%%% Merged introduction outline for "Proper Calibeating" — by Claude
%%% Hybrid of Claude/GPT/Gemini outlines.
%%%
%%% 7-section structure:
%%%   1. Setup and the question (+ thesis sentence up front)
%%%   2. Why this is nontrivial (stolen from GPT)
%%%   3. Calibration transfers for free
%%%   4. Calibeating does NOT transfer
%%%   5. Three routes to proper calibeating (+ joint-binning insight)
%%%   6. Comparison with concurrent work
%%%   7. Literature and roadmap

%% ====================================================================
%% 1. SETUP AND THE QUESTION
%% ====================================================================
%%
%% Brisk opening: proper scoring rules, Brier decomposition, calibeating.
%% The thesis sentence (from GPT) goes here so the reader knows
%% where the paper lands before wading into details.

\subsubsection*{1. Setup and the question}

A forecaster issues probability distributions; nature reveals outcomes;
the forecaster is evaluated by a scoring rule $L$.  A scoring rule is
\emph{proper} if truthful reporting minimizes expected loss; the
standard choice is the quadratic (Brier) scoring rule.

The Brier score decomposes as $\mathcal{B} = \mathcal{K} + \mathcal{R}$
(calibration plus refinement).  Calibration asks that forecasts match
realized frequencies; calibeating asks for a Brier score no worse than
the refinement of a reference forecaster---thus gaining calibration
without losing expertise.

All existing guarantees---stochastic calibration (Foster and Vohra
1998), continuous calibration (Foster and Hart 2021), and calibeating
(Foster and Hart 2023)---are stated exclusively for the Brier score.
But there are infinitely many proper scoring rules.  A natural
robustness question:

\medskip
\begin{quote}
\emph{If a procedure is calibrated (or calibeats) under the Brier
score, do these guarantees persist under every other proper scoring
rule?}
\end{quote}

\medskip\noindent
The answer, in one sentence:
\emph{Calibration is universally robust across proper Lipschitz scoring
rules, but calibeating is not; this paper identifies exactly what
additional structure restores that robustness.}

We work uniformly over $\mathcal{L}_{1}$, the class of all proper
$1$-Lipschitz scoring rules, and exclude rules with infinite losses
(such as the logarithmic rule).

\medskip\noindent\textbf{Cite:}
Brier (1950),
DeGroot and Fienberg (1983),
Foster and Vohra (1998),
Foster and Hart (2021),
Foster and Hart (2023).
%% Background on proper scoring rules: Savage (1971), McCarthy (1956),
%% Gneiting and Raftery (2007).  On calibration: Dawid (1982).


%% ====================================================================
%% 2. WHY THIS IS NONTRIVIAL
%% ====================================================================
%%
%% Stolen from GPT §2.  Earns the reader's attention before giving
%% the answer: the B=K+R decomposition is special to Brier; once you
%% leave Brier, that decomposition is not automatic.

\subsubsection*{2. Why this is nontrivial}

The Brier score is special.  The decomposition
$\mathcal{B} = \mathcal{K} + \mathcal{R}$ is an \emph{exact} identity
for the quadratic scoring rule: the Brier score equals calibration
error plus refinement, with nothing left over.  All known proofs of
calibeating exploit this decomposition.

For a general proper scoring rule $L$, an analogous decomposition
$\mathcal{B}^{L} = \mathcal{K}^{L} + \mathcal{R}^{L}$ holds when the
binning refines the forecast partition, but it breaks down for
fractional or continuous binning schemes---exactly the schemes used by
the strongest calibration procedures.  So the transfer of
guarantees from Brier to general proper rules is not automatic; it
requires new inequalities and, in some cases, new procedures.


%% ====================================================================
%% 3. CALIBRATION TRANSFERS FOR FREE
%% ====================================================================

\subsubsection*{3. Result: calibration is automatically proper-calibration}

Our first result is good news: calibration transfers for free across all
proper Lipschitz scoring rules.

Any $\varepsilon$-calibrated procedure is automatically
$\varepsilon$-\emph{proper}-calibrated.  The key inequality is
\[
\mathcal{K}^{L}_{t}(\mathbf{c};\mathbf{f})
\;\leq\; M\,\mathcal{K}_{t}(\mathbf{c};\mathbf{f})
\]
for every proper $M$-Lipschitz scoring rule $L$: the $L$-calibration
score is bounded by the Lipschitz constant times the quadratic
calibration score.  This holds for stochastic procedures (Foster and
Vohra 1998 type), for almost-deterministic procedures, and for
deterministic continuously calibrated procedures (Foster and Hart 2021
type).

In particular, all known calibrated forecasting procedures are already
proper-calibrated---no modification needed.

\medskip\noindent\textbf{Cite:}
Foster and Vohra (1998),
Foster and Hart (2021),
Foster (1999),
Blackwell (1956),  %% approachability — engine behind stochastic calibration
Dawid (1982),
Kakade and Foster (2008),
Hart (2022),
Mannor and Stoltz (2010).


%% ====================================================================
%% 4. CALIBEATING DOES NOT TRANSFER
%% ====================================================================

\subsubsection*{4. Result: calibeating does NOT imply proper calibeating}

The second result is a surprise: unlike calibration, calibeating does
\emph{not} transfer.

We exhibit a concrete example where a forecasting sequence $\mathbf{c}$
calibeats a reference $\mathbf{b}$ under the quadratic scoring rule but
\emph{fails} to calibeat $\mathbf{b}$ under the $2$-spherical scoring
rule (which is also proper and Lipschitz).  Both sequences are perfectly
calibrated; the failure is purely in the refinement comparison.
Repeating the example periodically gives the gap in the limit.

The reason calibration transfers but calibeating does not: calibration
involves only $\mathcal{K}^{L}$, which is controlled by
$\mathcal{K}$ via the bound above.  Calibeating involves
$\mathcal{B}^{L} - \mathcal{R}^{L}$, and there is no analogous
inequality relating $\mathcal{R}^{L}$ to $\mathcal{R}$.

\medskip\noindent\textbf{Cite:}
Foster and Hart (2023),
Gneiting and Raftery (2007).  %% spherical scoring rule family


%% ====================================================================
%% 5. THREE ROUTES TO PROPER CALIBEATING
%% ====================================================================
%%
%% Merged from Claude §4.  The joint-binning insight (stolen from
%% GPT §7) is promoted to a paragraph between (a) and (b).

\subsubsection*{5. Results: proper calibeating}

We provide three ways to achieve proper calibeating, plus
multicalibeating.

\medskip\noindent\textbf{(a) The simple procedure works.}
The deterministic procedure of Theorem~3 of Foster and Hart (2023),
which forecasts the running average action in the reference forecaster's
bin ($c_{t} = \bar{a}_{t-1}^{\mathbf{b}}(b_{t})$), is in fact
proper-calibeating.  The proof uses a new bound: online refinement
$\widetilde{\mathcal{R}}^{L}$ converges to offline refinement
$\mathcal{R}^{L}$ at rate $O((\ln t)/t)$ for every proper
$M$-Lipschitz scoring rule $L$.  This generalizes the quadratic-case
result of Foster and Hart (2023, Proposition~2).

\medskip\noindent\textbf{Why joint binning matters.}
The obstruction in Section~4 is not merely changing the scoring rule; it
is losing the refinement structure when bins mix different forecasts.
Calibeating plus calibration is \emph{not} sufficient for proper
calibeating.  What restores the decomposition logic is calibeating the
\emph{joint} binning $\mathbf{b}\times\mathbf{c}$, because the joint
always refines the forecast partition.  The counterexample of Section~4
illustrates exactly this gap.

\medskip\noindent\textbf{(b) Calibeating the joint binning.}
A stochastic procedure that calibeats the joint
$\mathbf{b}\times\mathbf{c}$ binning is both $(\delta,B)$-proper-%
calibeating and $\delta$-proper-calibrated.  The argument chains
the decomposition $\mathcal{B}^{L} = \mathcal{K}^{L}(\mathbf{c};
\mathbf{b}\times\mathbf{c}) + \mathcal{R}^{L}(\mathbf{b}\times
\mathbf{c})$ with the bound $\mathcal{K}^{L} \leq M\,\mathcal{K}$,
reducing proper calibeating of the joint to ordinary calibeating of the
joint (Theorem~5 of FH2023).

\medskip\noindent\textbf{(c) Deterministic continuous version.}
A deterministic continuously proper-calibrated procedure that
proper-calibeats.  Since fractional binnings do not refine the
forecast-binning, the exact decomposition
$\mathcal{B}^{L} = \mathcal{K}^{L} + \mathcal{R}^{L}$ breaks down.
We replace it with an approximate version:
$|\mathcal{B}^{L} - (\mathcal{K}^{L} + \mathcal{R}^{L})| < 2M\delta$
for any $\delta$-local fractional binning under an $M$-Lipschitz rule.
This generalizes the quadratic approximate decomposition from the errata
(Foster and Hart 2026), and is the key technical tool for this result.

\medskip\noindent\textbf{(d) Multicalibeating.}
All three results extend to simultaneously proper-calibeating
$N$ reference forecasters, by applying (a)--(c) to the product binning
$\mathbf{b}^{1}\times\cdots\times\mathbf{b}^{N}$.

\medskip\noindent\textbf{Cite:}
Foster and Hart (2023)---Theorems 3, 5, 7 being generalized,
Foster and Hart (2021)---continuous calibration framework,
Foster and Hart (2026)---errata / approximate decomposition,
Foster and Vohra (1998)---stochastic calibration machinery,
DeGroot and Fienberg (1983)---the original $\mathcal{B}=\mathcal{K}+
\mathcal{R}$ decomposition,
Lee, Noarov, Pai, and Roth (2022)---multicalibeating.


%% ====================================================================
%% 6. COMPARISON WITH CONCURRENT AND RELATED WORK
%% ====================================================================

\subsubsection*{6. Comparison with concurrent and related work}

Three concurrent lines of work address universality across scoring
rules, each from a different angle.  A careful comparison:

\medskip\noindent\textbf{Chen, Huang, Jordan, and Luo (2026).}
``Calibeating Made Simple'' reduces calibeating to regret minimization
and obtains calibeating rates for general proper losses, including
Brier, log, and mixable losses.  Their approach works loss by loss: for
each proper loss $L$, they build a procedure achieving
$\mathcal{B}^{L} \leq \mathcal{R}^{L} + O((\log T)/T)$.  Our
approach is different: we show that a \emph{single} procedure (the
existing Brier-score procedures) achieves proper calibeating
\emph{simultaneously} for all $L \in \mathcal{L}_{1}$.  The uniformity
is a consequence of the $\mathcal{K}^{L} \leq M\,\mathcal{K}$ bound,
which has no analogue for refinement.  They also handle the logarithmic
rule (which is not Lipschitz, so outside our class).

\medskip\noindent\textbf{Kleinberg, Leme, Schneider, and Teng (2023).}
``U-Calibration'' asks: what forecasting criterion guarantees low regret
for an unknown agent who may use any bounded scoring rule?  They show
calibration suffices but is not necessary, and introduce U-calibration
as the tight criterion, achieving $O(\sqrt{T})$ U-calibration error.
Our work is complementary: we are not introducing a new criterion but
showing that existing calibration/calibeating procedures already provide
uniform guarantees over all proper Lipschitz rules.

\medskip\noindent\textbf{Gopalan, Kalai, Reingold, Sharan, and Wieder
(2022).}  Omnipredictors show that multicalibration implies good
performance for all convex Lipschitz loss functions.  Our universality
result is analogous but in a different setting: we show calibration
implies universality over \emph{evaluation rules} (scoring rules), not
over \emph{decision-making losses}.  The mechanism is also different
(the $\mathcal{K}^{L} \leq M\,\mathcal{K}$ bound, not
omniprediction).

\medskip\noindent\textbf{Luo, Senapati, and Sharan (2025).}
Simultaneous swap regret minimization via KL-calibration achieves
$O(T^{1/3})$ swap regret for all proper losses with smooth univariate
form.  Another route to universality, via the swap-regret/calibration
connection rather than the calibeating/refinement decomposition.

\medskip\noindent\textbf{Dimitriadis, Gneiting, and Jordan (2023).}
Decompose the CRPS into miscalibration, discrimination, and uncertainty
components, relating to a hierarchy of calibration notions.  Study the
calibration--refinement decomposition for one specific proper rule (CRPS)
from a statistical perspective; we provide the game-theoretic counterpart.

\medskip\noindent\textbf{Popordanoska, Uria, and Blundell (2023).}
Show that every proper score decomposes into proper calibration error plus
refinement via a Bregman divergence; provide consistent estimators.
Complementary to our work: they estimate the decomposition, we exploit
it for uniform guarantees.

\medskip\noindent\textbf{Blasiok, Gopalan, Hu, and Nakkiran (2023).}
The converse direction: local optimality of a proper loss implies smooth
calibration.  We go in the other direction (from calibration to
universality over scoring rules).


%% ====================================================================
%% 7. LITERATURE AND ROADMAP
%% ====================================================================

\subsubsection*{7. Literature and roadmap}

\noindent\textbf{Proper scoring rules:}
Brier (1950),
McCarthy (1956),
Savage (1971),
Schervish (1989),
Gneiting and Raftery (2007),
Gneiting, Balabdaoui, and Raftery (2007),
Jose, Nau, and Winkler (2008),
Waghmare and Ziegel (2025).

\medskip\noindent\textbf{Calibration:}
Dawid (1982),
Foster and Vohra (1998),
Foster (1999),
Blackwell (1956),
Hart and Mas-Colell (2000),
Kakade and Foster (2008),
Mannor and Stoltz (2010),
Perchet (2011),
Rakhlin, Sridharan, and Tewari (2011),
Foster and Hart (2018),
Hart (2022),
Olszewski (2015).

\medskip\noindent\textbf{Calibeating:}
Foster and Hart (2023, 2022, 2024, 2026),
Lee, Noarov, Pai, and Roth (2022),
Gupta and Ramdas (2023),
Jain and Perchet (2024).

\medskip\noindent\textbf{Multicalibration and omniprediction:}
H\'{e}bert-Johnson, Kim, Reingold, and Rothblum (2018),
Gupta, Jung, Noarov, Pai, and Roth (2021),
Gopalan, Kalai, Reingold, Sharan, and Wieder (2022),
Hu, Luo, Senapati, and Sharan (2025),
Kleinberg, Mullainathan, and Raghavan (2016).

\medskip\noindent\textbf{Score decompositions:}
Dimitriadis, Gneiting, and Jordan (2023),
Popordanoska, Uria, and Blundell (2023).

\medskip\noindent\textbf{Other:}
Vovk (2009),
Qiao and Valiant (2021),
Marx, Kuleshov, and Ermon (2024),
Okoroafor, Kleinberg, and Sun (2023).

\medskip\noindent
Section~2 sets up the framework.  Sections~3--4 develop the proper
decomposition.  Section~5 proves the calibration transfer.  Section~6
gives the counterexample.  Section~7 establishes the three routes to
proper calibeating.  Section~8 concludes.
